```latex
\section{A low-order dynamical model for chaos}

The analytical solutions developed in this work describe tornado-like vortices under prescribed forcing and therefore do not address the question of how such vortices emerge dynamically. An alternative viewpoint is to seek a reduced-order dynamical system that captures the essential nonlinear interactions leading to tornado formation. Such a model would not attempt to reproduce the full complexity of atmospheric convection, but rather identify the minimal physical ingredients responsible for vortex genesis, maintenance, and possible chaotic evolution.

Motivated by the axisymmetric Navier--Stokes equations, we consider three variables representing the dominant dynamical processes,

\begin{equation}
x(t) = \text{radial inflow},
\qquad
y(t) = \text{azimuthal swirl},
\qquad
z(t) = \text{vertical stretching (updraft)}.
\end{equation}

These variables correspond to the three principal mechanisms governing tornado dynamics: low-level convergence, vortex rotation, and vertical stretching.

The first coupling follows directly from incompressibility. For an axisymmetric flow,

\begin{equation}
\frac{\partial u_r}{\partial r}
+
\frac{u_r}{r}
+
\frac{\partial w}{\partial z}
=
0,
\end{equation}

so that stronger radial convergence necessarily enhances the upward velocity. Consequently, radial inflow acts as a source for vertical stretching,

\begin{equation}
x \longrightarrow z.
\end{equation}

The second coupling arises from conservation of angular momentum. Neglecting viscous diffusion, the azimuthal momentum equation satisfies

\begin{equation}
\frac{D}{Dt}(r v_\theta)
\approx
0,
\end{equation}

implying that vertical stretching decreases the effective vortex radius and thereby amplifies the azimuthal velocity. This positive feedback may be represented schematically as

\begin{equation}
z \longrightarrow y.
\end{equation}

Finally, the radial momentum equation contains the centrifugal acceleration

\begin{equation}
-\frac{v_\theta^2}{r},
\end{equation}

which opposes radial convergence. As the swirl intensity increases, the resulting centrifugal force weakens the inflow, giving the negative feedback

\begin{equation}
y \longrightarrow -x.
\end{equation}

Together these three processes form the nonlinear feedback cycle

\begin{equation}
x
\rightarrow
z
\rightarrow
y
\rightarrow
-x,
\end{equation}

which resembles the cyclic interactions found in the Lorenz equations and other low-dimensional chaotic systems.

A phenomenological model consistent with these physical mechanisms may therefore be written as

\begin{subequations}
\begin{align}
\dot{x}
&=
-\sigma x
-\alpha y^2
+\beta z
+F,
\\
\dot{y}
&=
-\mu y
+\gamma xz,
\\
\dot{z}
&=
-r z
+\delta x
-\eta y,
\end{align}
\end{subequations}

where $F$ represents the environmental forcing, while the remaining coefficients parameterize dissipation and nonlinear coupling strengths.

The physical interpretation of each term is straightforward. The linear terms describe frictional decay and turbulent dissipation. The forcing $F$ represents the large-scale environment that drives low-level convergence. The quadratic term $-\alpha y^2$ models the centrifugal inhibition of radial inflow by strong vortex rotation. The nonlinear production term $\gamma xz$ represents the amplification of swirl through the combined effects of convergence and vertical stretching. Finally, the terms $\beta z$ and $\delta x$ express the mutual coupling between convergence and the updraft through mass continuity.

Within this framework, tornado formation may be interpreted as a nonlinear bifurcation. For weak environmental forcing, the trivial equilibrium

\begin{equation}
(x,y,z)
=
(0,0,0)
\end{equation}

represents a quiescent atmosphere without organized rotation. As the forcing increases, a finite-amplitude equilibrium may emerge, corresponding to a sustained tornado-like vortex. Further parameter changes may destabilize this equilibrium through Hopf or period-doubling bifurcations, producing oscillatory or chaotic dynamics. Such transitions provide a possible dynamical explanation for the observed variability of tornado intensity, cyclic evolution, and rapid structural changes.

Although the above system is introduced phenomenologically, it provides a useful conceptual framework for understanding tornado dynamics. A more rigorous reduced-order model may be obtained by deriving evolution equations directly from exact or approximate Navier--Stokes solutions. For example, if an analytical vortex solution is parameterized by quantities such as circulation, strain rate, and vertical shear, these parameters may be promoted to time-dependent variables whose evolution is determined through projection or modulation methods. Such an approach would retain a direct connection with the governing equations while preserving the analytical vortex structure, thereby providing a mathematically consistent bridge between exact solutions of the Navier--Stokes equations and low-dimensional dynamical systems capable of exhibiting bifurcations and chaotic behavior.
```
